% GENERATED FILE. Edit canonical lessons, then run npm run generate:books.
% canonical-source-hash: fnv1a64:62aaaa290549c7ec
% canonical-lessons: ML-C519-utkantha, ML-C519-atmabhimanam, ML-C519-amparakkuka, ML-C519-vatsalyam, ML-C519-murumurukkuka

\chapter{Worry and Relief}
\label{ch:ml-worry-and-relief}
\hlchaptermodality{519}

\begin{chapteropening}
\textbf{By the end of this chapter:} \emph{I can say anxiety, self-respect, to be astonished, tender affection and to grumble.}

Complete the last lesson of chapter 519: I can say anxiety, self-respect, to be astonished, tender affection and to grumble.
\end{chapteropening}

\section[\texorpdfstring{\ml{ഉത്കണ്ഠ} (utkaṇṭha)}{ഉത്കണ്ഠ (utkaṇṭha)}]{\ml{ഉത്കണ്ഠ} (utkaṇṭha) — anxiety; nervous worry}

\label{lesson:ML-C519-utkantha}



\begin{quote}
Before the new one: say the Malayalam for to long for, then the Malayalam for bliss.
\end{quote}

\subsection*{You'll want to know: \ml{ഉത്കണ്ഠ}}
\textbf{\ml{ഉത്കണ്ഠ}} — \emph{utkaṇṭha} — \textquotedblleft{}anxiety; nervous worry\textquotedblright{}.

\begin{quote}
\ml{പരീക്ഷയ്ക്ക്} \ml{മുൻപ്} \ml{അവൾക്ക്} \ml{വലിയ} \ml{ഉത്കണ്ഠ} \ml{ഉണ്ടായിരുന്നു}.
\end{quote}

\begin{quote}
\emph{parīkṣaykkŭ munpŭ avaḷkkŭ valiya utkaṇṭha uṇṭāyirunnu.}
\end{quote}

\begin{quote}
She was very anxious before the exam.
\end{quote}

Formal and close to \ml{ആശങ്ക}; the one who feels it takes -\ml{ക്ക്}.

\begin{cousinweb}
\textbf{\ml{ഉത്കണ്ഠ}} (\emph{utkaṇṭha}) is borrowed from Sanskrit \emph{utkaṇṭhā}, \textquotedblleft{}longing for someone, missing them\textquotedblright{} (Monier-Williams). In Malayalam the anxious waiting in that longing became the main sense.
\end{cousinweb}

\begin{grammarlens}[title={What you've built}]
One more word for describing what you lived and felt.
\end{grammarlens}

\subsection*{Your turn}
\begin{itemize}
\raggedright
  \item \emph{Say it:} \emph{utkaṇṭha}
  \item \emph{Say it:} \emph{utkaṇṭha}, once more
  \item \emph{Say it:} \emph{parīkṣaykkŭ munpŭ avaḷkkŭ valiya utkaṇṭha uṇṭāyirunnu.} — She was very anxious before the exam.
  \item \emph{Recall:} say the Malayalam for to long for, then the Malayalam for bliss, then say \emph{utkaṇṭha} again
\end{itemize}

\begin{tcolorbox}[breakable,colback=teal!4,colframe=teal!35!black,title={Before you move on}]
What does \ml{ഉത്കണ്ഠ} mean? (\textquotedblleft{}Anxiety; nervous worry\textquotedblright{}.) Say it once more.
\end{tcolorbox}
\section[\texorpdfstring{\ml{ആത്മാഭിമാനം} (ātmābhimānaṁ)}{ആത്മാഭിമാനം (ātmābhimānaṁ)}]{\ml{ആത്മാഭിമാനം} (ātmābhimānaṁ) — self-respect}

\label{lesson:ML-C519-atmabhimanam}



\begin{quote}
Before the new one: say the Malayalam for bliss, then the Malayalam for anxiety.
\end{quote}

\subsection*{You'll want to know: \ml{ആത്മാഭിമാനം}}
\textbf{\ml{ആത്മാഭിമാനം}} — \emph{ātmābhimānaṁ} — \textquotedblleft{}self-respect\textquotedblright{}.

\begin{quote}
\ml{ആ} \ml{വാക്കുകൾ} \ml{അവന്റെ} \ml{ആത്മാഭിമാനത്തെ} \ml{വേദനിപ്പിച്ചു}.
\end{quote}

\begin{quote}
\emph{ā vākkukaḷ avanṟe ātmābhimānatte vēdanippiccu.}
\end{quote}

\begin{quote}
Those words hurt his self-respect.
\end{quote}

\ml{ആത്മ}- self + \ml{അഭിമാനം} pride: pride in oneself, already familiar as \ml{അഭിമാനം}.

\begin{cousinweb}
\textbf{\ml{ആത്മാഭിമാനം}} (\emph{ātmābhimānaṁ}) is a learned compound built from Sanskrit \emph{ātman}, \textquotedblleft{}self\textquotedblright{} (Monier-Williams), and \emph{abhimāna}, \textquotedblleft{}pride\textquotedblright{} (CDIAL 524).
\end{cousinweb}

\begin{grammarlens}[title={What you've built}]
One more word for describing what you lived and felt.
\end{grammarlens}

\subsection*{Your turn}
\begin{itemize}
\raggedright
  \item \emph{Say it:} \emph{ātmābhimānaṁ}
  \item \emph{Say it:} \emph{ātmābhimānaṁ}, once more
  \item \emph{Say it:} \emph{ā vākkukaḷ avanṟe ātmābhimānatte vēdanippiccu.} — Those words hurt his self-respect.
  \item \emph{Recall:} say the Malayalam for bliss, then the Malayalam for anxiety, then say \emph{ātmābhimānaṁ} again
\end{itemize}

\begin{tcolorbox}[breakable,colback=teal!4,colframe=teal!35!black,title={Before you move on}]
What does \ml{ആത്മാഭിമാനം} mean? (\textquotedblleft{}Self-respect\textquotedblright{}.) Say it once more.
\end{tcolorbox}
\section[\texorpdfstring{\ml{അമ്പരക്കുക} (amparakkuka)}{അമ്പരക്കുക (amparakkuka)}]{\ml{അമ്പരക്കുക} (amparakkuka) — to be astonished; to be bewildered}

\label{lesson:ML-C519-amparakkuka}



\begin{quote}
Before the new one: say the Malayalam for anxiety, then the Malayalam for self-respect.
\end{quote}

\subsection*{You'll want to know: \ml{അമ്പരക്കുക}}
\textbf{\ml{അമ്പരക്കുക}} — \emph{amparakkuka} — \textquotedblleft{}to be astonished; to be bewildered\textquotedblright{}.

\begin{quote}
\ml{വലിയ} \ml{നഗരം} \ml{കണ്ട്} \ml{ഗ്രാമത്തിലെ} \ml{കുട്ടി} \ml{അമ്പരന്നു}.
\end{quote}

\begin{quote}
\emph{valiya nagaraṁ kaṇṭŭ grāmattile kuṭṭi amparannu.}
\end{quote}

\begin{quote}
The village child was bewildered by the big city.
\end{quote}

Past: \ml{അമ്പരന്നു}, \emph{amparannu}.

The past is \ml{അമ്പരന്നു}, with -\ml{ന്നു}. The noun is \ml{അമ്പരപ്പ്}, astonishment.

\begin{cousinweb}
\textbf{\ml{അമ്പരക്കുക}} (\emph{amparakkuka}) is native Dravidian: DEDR 172 gives the Malayalam verb as to be confounded or perplexed, with the noun \ml{അമ്പരപ്പ്}, perplexity (DEDR 172).
\end{cousinweb}

\begin{grammarlens}[title={What you've built}]
One more word for describing what you lived and felt.
\end{grammarlens}

\subsection*{Your turn}
\begin{itemize}
\raggedright
  \item \emph{Say it:} \emph{amparakkuka}
  \item \emph{Say it:} \emph{amparakkuka}, once more
  \item \emph{Say it:} \emph{valiya nagaraṁ kaṇṭŭ grāmattile kuṭṭi amparannu.} — The village child was bewildered by the big city.
  \item \emph{Recall:} say the Malayalam for anxiety, then the Malayalam for self-respect, then say \emph{amparakkuka} again
\end{itemize}

\begin{tcolorbox}[breakable,colback=teal!4,colframe=teal!35!black,title={Before you move on}]
What does \ml{അമ്പരക്കുക} mean? (\textquotedblleft{}To be astonished; to be bewildered\textquotedblright{}.) Say it once more.
\end{tcolorbox}
\section[\texorpdfstring{\ml{വാത്സല്യം} (vātsalyaṁ)}{വാത്സല്യം (vātsalyaṁ)}]{\ml{വാത്സല്യം} (vātsalyaṁ) — tender affection}

\label{lesson:ML-C519-vatsalyam}



\begin{quote}
Before the new one: say the Malayalam for self-respect, then the Malayalam for to be astonished.
\end{quote}

\subsection*{You'll want to know: \ml{വാത്സല്യം}}
\textbf{\ml{വാത്സല്യം}} — \emph{vātsalyaṁ} — \textquotedblleft{}tender affection\textquotedblright{}.

\begin{quote}
\ml{മുത്തശ്ശിക്ക്} \ml{പേരക്കുട്ടികളോട്} \ml{വലിയ} \ml{വാത്സല്യമാണ്}.
\end{quote}

\begin{quote}
\emph{muttaśśikkŭ pērakkuṭṭikaḷōṭŭ valiya vātsalyamāṇŭ.}
\end{quote}

\begin{quote}
Grandmother has great affection for her grandchildren.
\end{quote}

The tender love of an older person for a younger one; the one loved takes -\ml{ഓട്}.

\begin{cousinweb}
\textbf{\ml{വാത്സല്യം}} (\emph{vātsalyaṁ}) is borrowed from Sanskrit \emph{vātsalya}, \textquotedblleft{}affection\textquotedblright{} (CDIAL 11506), the noun of \emph{vatsala}, \textquotedblleft{}fond, tender\textquotedblright{}, a word built on \emph{vatsa}, \textquotedblleft{}a calf\textquotedblright{}: the love of a cow for her calf.
\end{cousinweb}

\begin{grammarlens}[title={What you've built}]
One more word for describing what you lived and felt.
\end{grammarlens}

\subsection*{Your turn}
\begin{itemize}
\raggedright
  \item \emph{Say it:} \emph{vātsalyaṁ}
  \item \emph{Say it:} \emph{vātsalyaṁ}, once more
  \item \emph{Say it:} \emph{muttaśśikkŭ pērakkuṭṭikaḷōṭŭ valiya vātsalyamāṇŭ.} — Grandmother has great affection for her grandchildren.
  \item \emph{Recall:} say the Malayalam for self-respect, then the Malayalam for to be astonished, then say \emph{vātsalyaṁ} again
\end{itemize}

\begin{tcolorbox}[breakable,colback=teal!4,colframe=teal!35!black,title={Before you move on}]
What does \ml{വാത്സല്യം} mean? (\textquotedblleft{}Tender affection\textquotedblright{}.) Say it once more.
\end{tcolorbox}
\section[\texorpdfstring{\ml{മുറുമുറുക്കുക} (muṟumuṟukkuka)}{മുറുമുറുക്കുക (muṟumuṟukkuka)}]{\ml{മുറുമുറുക്കുക} (muṟumuṟukkuka) — to grumble}

\label{lesson:ML-C519-murumurukkuka}



\begin{quote}
Before the new one: say the Malayalam for to be astonished, then the Malayalam for tender affection.
\end{quote}

\subsection*{You'll want to know: \ml{മുറുമുറുക്കുക}}
\textbf{\ml{മുറുമുറുക്കുക}} — \emph{muṟumuṟukkuka} — \textquotedblleft{}to grumble\textquotedblright{}.

\begin{quote}
\ml{കൂടുതൽ} \ml{ജോലി} \ml{കിട്ടിയപ്പോൾ} \ml{ജീവനക്കാർ} \ml{മുറുമുറുത്തു}.
\end{quote}

\begin{quote}
\emph{kūṭutal jōli kiṭṭiyappōḷ jīvanakkār muṟumuṟuttu.}
\end{quote}

\begin{quote}
The staff grumbled when they were given more work.
\end{quote}

Past: \ml{മുറുമുറുത്തു}, \emph{muṟumuṟuttu}.

Complaining under one's breath, not out loud; the past is \ml{മുറുമുറുത്തു}.

\begin{cousinweb}
\textbf{\ml{മുറുമുറുക്കുക}} (\emph{muṟumuṟukkuka}) is native Dravidian: DEDR 5013 gives the Malayalam verb as to murmur (DEDR 5013). The doubled sound copies the low noise of a grumble.
\end{cousinweb}

\begin{grammarlens}[title={What you've built}]
One more word for describing what you lived and felt.
\end{grammarlens}

\subsection*{Your turn}
\begin{itemize}
\raggedright
  \item \emph{Say it:} \emph{muṟumuṟukkuka}
  \item \emph{Say it:} \emph{muṟumuṟukkuka}, once more
  \item \emph{Say it:} \emph{kūṭutal jōli kiṭṭiyappōḷ jīvanakkār muṟumuṟuttu.} — The staff grumbled when they were given more work.
  \item \emph{Recall:} say the Malayalam for to be astonished, then the Malayalam for tender affection, then say \emph{muṟumuṟukkuka} again
\end{itemize}

\begin{tcolorbox}[breakable,colback=teal!4,colframe=teal!35!black,title={Before you move on}]
What does \ml{മുറുമുറുക്കുക} mean? (\textquotedblleft{}To grumble\textquotedblright{}.) Say it once more.
\end{tcolorbox}
